\section{ROHF orbital energies are not unique (canonicalisation)}\label{sec:rohf}

\begin{reviewclaim}
``\dots the orbitals themselves are neither observable, nor unique (in fact, erroneous non-aufbau
occupations can be observed with restricted open-shell methods depending on the canonicalization
parameters of the $\alpha$ and $\beta$ electrons in Fock matrix, see Ref 82) \dots'' (p.~449; ref.~82
Glaesemann \& Schmidt 2010)
\end{reviewclaim}

\subsection{High-spin ROHF and its stationarity conditions}
Orbital spaces: doubly occupied (core) $c$, singly occupied ($\alpha$ only) $o$, virtual $v$. Densities:
$\mathbf D^\alpha=\mathbf D_c+\mathbf D_o$ and $\mathbf D^\beta=\mathbf D_c$. The energy has the UHF form
$E[\mathbf D^\alpha,\mathbf D^\beta]$, so its first variation is
\begin{equation}\label{eq:dE}
  \delta E=\operatorname{tr}(\mathbf F^\alpha\,\delta\mathbf D^\alpha)+\operatorname{tr}(\mathbf F^\beta\,\delta\mathbf D^\beta),
\end{equation}
with $\mathbf F^\sigma=\partial E/\partial\mathbf D^\sigma$ the UHF Fock matrices built from the ROHF densities.
\begin{enumerate}
  \item An infinitesimal real rotation mixing orbital $q$ into orbital $p$ by angle $\kappa$ changes the density of a
        spin $\sigma$ only if exactly one of $p,q$ is occupied in that spin. Then
        $\delta\mathbf D^\sigma=\kappa(\ket p\bra q+\ket q\bra p)$, and $\delta E$ picks up $2\kappa F^\sigma_{pq}$.
  \item \textbf{$c$--$o$ rotation:} for $\alpha$ both orbitals are occupied (no change); for $\beta$, $c$ is occupied
        and $o$ empty. Stationarity requires $F^\beta_{co}=0$.
  \item \textbf{$o$--$v$ rotation:} for $\alpha$, $o$ is occupied and $v$ empty; for $\beta$ both are empty.
        Stationarity requires $F^\alpha_{ov}=0$.
  \item \textbf{$c$--$v$ rotation:} both spins change. Stationarity requires $F^\alpha_{cv}+F^\beta_{cv}=0$.
  \item \textbf{Rotations within $c$, within $o$, or within $v$} leave $\mathbf D^\alpha$ and $\mathbf D^\beta$ unchanged, hence
        $\delta E=0$ \emph{identically}. Nothing in the variational principle fixes them.
\end{enumerate}

\subsection{Effective Fock operator and the arbitrariness of orbital energies}
To get orbitals from a single eigenvalue problem, one builds an effective Fock matrix $\mathbf R$ whose
off-diagonal blocks enforce steps 2--4:
\begin{equation}
  \mathbf R=\begin{pmatrix}
  \mathbf R_{cc} & \mathbf F^\beta_{co} & \tfrac12(\mathbf F^\alpha+\mathbf F^\beta)_{cv}\\
  \cdot & \mathbf R_{oo} & \mathbf F^\alpha_{ov}\\
  \cdot & \cdot & \mathbf R_{vv}\end{pmatrix},\qquad
  \mathbf R_{xx}=A_{xx}\mathbf F^\alpha_{xx}+B_{xx}\mathbf F^\beta_{xx}\ \ (\text{any choice}).
\end{equation}
At convergence the off-diagonal blocks vanish (steps 2--4), and the orbitals diagonalise each diagonal
block separately. By step 5 the total energy, the densities and every observable derived from them are
independent of $A_{xx},B_{xx}$. The \emph{orbital energies}, the eigenvalues of $\mathbf R_{xx}$, are not. The
relative position of the SOMO (eigenvalue of $\mathbf R_{oo}$) and the HOMO of the closed shell (top
eigenvalue of $\mathbf R_{cc}$) is therefore a matter of convention.

\begin{verifybox}
T11a (\texttt{t11\_rohf\_and\_distonic.py}), ROHF/6-31G*:
\begin{itemize}
  \item Stationarity: $|F^\beta_{co}|,|F^\alpha_{ov}|,|(F^\alpha+F^\beta)_{cv}|<10^{-7}$ for H$_2$NO, NO$_2$, CH$_3$ and allyl.
        The blocks that need not vanish do not: $|F^\alpha_{co}|=0.11$, $|F^\beta_{ov}|=0.16$ (H$_2$NO). This confirms
        steps 2--4. PySCF's \texttt{get\_roothaan\_fock} uses exactly these off-diagonal blocks.
  \item Canonicalisation, H$_2$NO: $\varepsilon_{\text{SOMO}}=-0.119$ ($\mathbf R_{oo}=\tfrac12(\mathbf F^\alpha+\mathbf F^\beta)$),
        $-0.440$ ($\mathbf F^\alpha$), $+0.203\,E_h$ ($\mathbf F^\beta$), while the closed-shell HOMO is $-0.467$. Same energy
        $E=-130.36970936\,E_h$ in every case.
  \item With one open orbital, $\varepsilon_{\text{SOMO}}(A)=AF^\alpha_{oo}+(1-A)F^\beta_{oo}$ is linear in $A$. The SOMO drops
        below the HOMO for $A>A^*$, with $A^*\approx1.04$ for H$_2$NO and NO$_2$. The flip was verified explicitly.
\end{itemize}
\textbf{PASS}.
\end{verifybox}
\begin{caveatbox}
For these ordinary radicals the ordering flips only for weights slightly outside the ``convex'' range
$0\le A\le1$. So the review's warning is mathematically exact (the ordering is not invariant), but whether
a \emph{standard} canonicalisation produces a spurious non-aufbau ordering depends on the molecule and
the scheme. This is why the review, correctly, asks for evidence from observables: ionisation
energies, redox potentials and the spin state of the products (\cref{sec:qcs}).
\end{caveatbox}
\begin{sourcebox}
Glaesemann \& Schmidt, \emph{J.\ Phys.\ Chem.\ A} \textbf{114}, 8772 (2010), doi:10.1021/jp101758y: not accessed.
The specific canonicalisation parameter sets they compare could not be checked. The derivation above
is general.
\end{sourcebox}
